\scanpage{68}
\margin{\figure{\begin{tikzpicture}[x=.8cm,y=.8cm]
\draw[line width=.55pt] (0,0) .. controls (.12,-.02) and (.25,-.01) .. (.33,.02) -- (.29,.19);
\draw[line width=.75pt] (1.35,-.10) -- (1.64,-.13) -- (1.56,.15);
\end{tikzpicture}}}
\deleted{sous les conditions envisagées a) b) $\quad)$ $\quad)$, $B$ plat sur $A$ et à extension résiduelle finie,
\[
\text{\uuline{R}}(A)\quad\text{et}\quad\text{\uuline{R}}(B\otimes k)\ \text{géométrique}\quad\text{impliquent}\quad\text{\uuline{R}}(B).
\]
La condition supplémentaire qu'on vient de signaler pour \uuline{R} est satisfaite dans tous les exemples envisagés dans 7.5.3., à l'exception de (vii), où il est en défaut. En effet, on peut (même \deleted{sur un corps} pour une courbe $Y$ sur un corps algébriquement clos $k$ de caractéristique arbitraire) donner des exemples d'un morphisme étale $f\colon X\quad Y$, et d'un point $x$ de $X$ tel que $Y$ soit intègre en $y=f\deleted{x}(x)$, mais $X$ ne le soit pas en $x$.\handwritten{/} Ce serait peut-être le moment de donner l'exemple (également signalé dans 6.15.7. et qui resert $X$ fois): on prend une courbe à point double \deleted{\unsure{e|r|o}}o\deleted{\unsure{u|n}}rdinaire, et on recolle deux exemplaires du normalisé \unsure{d|de} de façon à avoir une courbe étale $X$ au dessus de $Y$. Un exemple analogue avec un anneau loc\deleted{\unsure{q|p}}al \inserted{intègre} \uline{complet}, \deleted{\unsure{a|u}}donc unibranche, mais en l'occurence non géométriquement unibranche, est donné par $\overset{\textstyle\inserted{A=}}{\text{\uuline{R}}}\ X,Y\quad/(X^{\handwritten{2}}+Y^{\handwritten{2}})$. \deleted{Si $X$ est son spectre} Alors $A\otimes_{\underline{R}}\underline{C}$ est étale sur $A$, mais non intègre. Ces exemples, qu'on semble rencontrer pour la première fois à propos de 6.15.7. pourraient peut-être remonter à cet endroit.\handwritten{/} Il est cependant p\deleted{o}lausible que \deleted{l'énoncé qui} la reponse à la question suivante soit affirmative :

Soit $A\quad B$ un homomorphisme local d'anneaux locaux noethériens, $B$ plat sur $A$, $B\otimes_A k\handwritten{/}k$ géométriquement intègre, $A$ \deleted{excellent (par exemple complet)} géométriquement unibranche (NB on peut se ramener au cas: $A$ normal). Si $A$ est de plus excellent (disons complet), alors $B$ est-il intègre ?

La reponse est affirmative quand le corps résiduel $k$ de $A$ est de car 0, comme on verra au N\textsuperscript{\uline{o}} 7.9. en utilisant la\deleted{é} résolution des singularités de Hironaka, ou lorsque $B$ est essentiellement de type fini sur $A$ (réf $\qquad$). M\unsure{\deleted{i}a\deleted{a}i|\deleted{a}i\deleted{i}a}s \begin{tabular}{@{}l@{}}\inserted{dans le cas général,}\\même lorsque l'extension\end{tabular} résiduelle est triviale, $A$\deleted{\unsure{e|s|x}} et $B$ complet, \deleted{la technique dé\unsure{vis|v\illegible{1}l\illegible{1}}} nous ignorons si la réponse est affirmative. On a cependant un résultat partiel ci-dessous.}
